\documentclass[12pt]{article}
\usepackage[T1]{fontenc}
\usepackage[utf8]{inputenc}
\usepackage{lmodern}
\usepackage{textcomp}
\usepackage{enumitem}
\usepackage{geometry}
\geometry{margin=1in}
\begin{document}
\begin{center}
Answer Key - Constitutional Law - Undergraduate Level Assessment 12 \
\small (Answers are ordered exactly as in the question paper.)
\end{center}

\section*{Multiple Choice}
\begin{enumerate}[label=\arabic*.]
\item \textbf{Answer:} D
\\ \textit{Options:} A) maximizing individual freedom; B) providing a basis for compromise; C) keeping the peace; D) promoting the principles of the free enterprise system
\\ \textit{Note:} The correct answer is "promoting the principles of the free enterprise system" because while the law does facilitate economic activities, its primary functions include maintaining order, resolving disputes, and protecting individual rights rather than explicitly promoting any specific economic system.
\item \textbf{Answer:} B
\\ \textit{Options:} A) interpret; B) make; C) enforce; D) override
\\ \textit{Note:} The legislative branch of the United States government, primarily composed of Congress, has the constitutional authority to create, amend, and enact laws, which is fundamental to its role in the American political system.
\item \textbf{Answer:} C
\\ \textit{Options:} A) The Depression of the 1930 s.; B) The rise of Fascism.; C) The international recognition of human rights after WWII.; D) The Bolshevik revolution.
\\ \textit{Note:} The international recognition of human rights after WWII contributed significantly to the revival of natural law in the 20 th century by re-establishing the idea that there are universal moral principles inherent to human dignity that should inform legal systems and governance.
\item \textbf{Answer:} D
\\ \textit{Options:} A) Conflict.; B) Love.; C) War.; D) Reciprocity.
\\ \textit{Note:} Malinowski identified reciprocity as a fundamental organizing principle of Trobriand society, emphasizing the importance of mutual exchange and social bonds in their social structure and economic interactions.
\item \textbf{Answer:} B
\\ \textit{Options:} A) Law is politics.; B) Law is determinate.; C) Law reflects economic power.; D) Law is unstable.
\\ \textit{Note:} Critical legal theorists reject the proposition that "law is determinate" because they argue that law is inherently indeterminate and influenced by social, political, and economic factors, making legal outcomes uncertain and subject to various interpretations.
\end{enumerate}

\section*{True / False}
\begin{enumerate}[label=\arabic*.]
\setcounter{enumi}{5}
\item \textbf{Answer:} True
\\ \textit{Options:} A) True; B) False
\\ \textit{Note:} H.L.A. Hart is indeed the author of 'The Concept of Law', published in 1961, which is a seminal work in legal philosophy that analyzes the structure and function of legal systems.
\item \textbf{Answer:} False
\\ \textit{Options:} A) True; B) False
\\ \textit{Note:} Austin's perspective on jurisprudence prioritizes positive law, asserting that law is a command issued by a sovereign and does not inherently involve moral principles.
\item \textbf{Answer:} False
\\ \textit{Options:} A) True; B) False
\\ \textit{Note:} The statement is false because many countries recognize the supremacy of international law over domestic law, particularly when they have ratified international treaties that obligate them to comply with those laws.
\item \textbf{Answer:} True
\\ \textit{Options:} A) True; B) False
\\ \textit{Note:} The International Covenant on Civil and Political Rights includes a reporting mechanism whereby states must submit regular reports on their compliance with the treaty, and it also allows individuals to submit petitions to the Human Rights Committee if they believe their rights under the covenant have been violated, thus confirming the statement as true.
\item \textbf{Answer:} False
\\ \textit{Options:} A) True; B) False
\\ \textit{Note:} Critical Race Theorists do not reject concepts like 'justice' but rather critique traditional understandings of justice as insufficient for addressing systemic racism and inequality.
\end{enumerate}

\section*{Long Form Response}
\begin{enumerate}[label=\arabic*.]
\setcounter{enumi}{10}
\item \textbf{Answer:} Jurisprudence, as articulated by John Austin, is understood as the philosophy of positive law, which emphasizes the importance of law as it is enacted by sovereign authority rather than moral considerations. Austin's legal positivism asserts that laws are commands backed by sanctions, and this view shapes constitutional law by providing a clear framework for interpreting legal texts based on their literal meaning and the intentions of the lawmakers. This approach has significant implications for legal interpretation, as it encourages judges and legal practitioners to focus on the application of statutes without delving into moral or ethical evaluations. By prioritizing the written law, Austin's theory supports a structured legal system that upholds the rule of law and facilitates predictable legal outcomes in constitutional contexts. Thus, understanding jurisprudence through Austin's lens is crucial for grasping the nature of legal authority and the interpretation of constitutional provisions.
\\ \textit{Note:} correct because it effectively captures Austin's emphasis on positive law as commands from a sovereign authority, highlighting its role in shaping constitutional law and guiding legal interpretation through a focus on the literal meaning of laws and the intentions behind them.
\item \textbf{Answer:} The entities entitled to request advisory opinions from the International Court of Justice (ICJ) include the United Nations General Assembly and the Security Council, which can seek guidance on any legal question. Additionally, other organs of the United Nations and specialized agencies may request advisory opinions if authorized by the General Assembly. These requests have significant implications for international law and governance, as they allow for the clarification and interpretation of legal principles that can influence state behavior and international relations. By providing advisory opinions, the ICJ contributes to the development of international law and helps ensure that legal norms are upheld in global governance. This process fosters a collaborative approach to addressing complex legal issues within the international community.
\\ \textit{Note:} correct because it accurately identifies the entities authorized to request advisory opinions from the ICJ, specifically the UN General Assembly and Security Council, and emphasizes the broader implications of these requests for the development and interpretation of international law and governance.
\end{enumerate}

\vfill
\noindent\textit{End of Answer Key}
\end{document}